% !TEX program = xelatex
\documentclass[12pt,a4paper]{article}
\usepackage{fontspec}
\setmainfont{Liberation Serif}
\usepackage[bulgarian]{babel}
\usepackage[a4paper,margin=2.5cm]{geometry}
\usepackage[expansion=false]{microtype}
\usepackage{hyperref}
\hypersetup{hidelinks,
  pdftitle={Резюме: Система за извличане с генериране за достъп до фармакотерапевтични ръководства на български език},
  pdfauthor={Цветелина Пашовска}}

\usepackage{indentfirst}
\setlength{\parindent}{1.25cm}
\setlength{\parskip}{0pt}
\widowpenalty=10000
\clubpenalty=10000
\pagestyle{empty}

\begin{document}

\begin{center}
{\small Софийски университет „Св. Климент Охридски“, Факултет по математика и информатика\\
Магистърска програма „Изкуствен интелект“}\\[0.8em]
{\large\bfseries Система за извличане с генериране (RAG) за достъп до\\ фармакотерапевтични ръководства на български език}\\[0.6em]
Дипломант: Цветелина Пашовска, фак. номер: 7MI3400669\\[0.8em]
Ръководител: проф. д-р Мария Нишева, катедра „Компютърна информатика“ \\[0.8em]
{\bfseries РЕЗЮМЕ}
\end{center}

Фармакотерапевтичните ръководства, утвърдени от Националния съвет по цени и реимбурсиране на лекарствените продукти, съдържат официални препоръки за лечение, но са обемни PDF файлове с разнородно оформление, част от тях сканирани. Анкета сред общопрактикуващи лекари потвърждава нуждата от по-лесен достъп до актуална медицинска информация, но и опасението от грешки. Не е открита публикувана система за извличане с генериране (RAG) за медицински ръководства на български език. Целта на работата е да се разработи прототип на такава система върху корпус от 1380 откъса от 17 национални ръководства, изграден в преддипломния проект.

Крайната система е разработена итеративно, като всяко решение е проверено с експеримент, а разликите между вариантите -- със статистически тестове. Извличането е хибридно: семантично търсене с модела bge-m3, избран след сравнение на три многоезични модела за векторно представяне, и лексикално търсене BM25 със съкращаване на думите, обединени чрез сливане по реципрочни рангове и подредени от модел за повторно класиране, чиято оценка служи и като праг за отказ. Отговорът се генерира локално от модела BgGPT с 4 милиарда параметри. Експериментите с подсказката показват, че малкият модел пренебрегва откъсите, отдалечени от въпроса, и измисля източници, когато бъде поискано да ги цитира. Затова най-релевантният откъс се поставя най-близо до въпроса, а източниците се показват от системата. Неподкрепените отговори се ограничават на няколко нива: проверка преди генерирането дали ключовите понятия от въпроса се срещат в откъсите, проверка на отговора за имена на лекарства, ATC кодове и числа, липсващи в откъсите, и маркиране на изреченията от източниците, на които стъпва отговорът. Уеб приложението работи изцяло локално на графичен процесор с 6 GB памет.

Окончателната оценка е проведена върху независим набор от 59 въпроса в обхвата на ръководствата и 10 медицински въпроса извън него, а отговорите са оценени ръчно с участието на общопрактикуващ лекар. Верният откъс е първи при 58\% от въпросите и е сред петте показани при 86\%, като повторното класиране значимо подобрява подреждането спрямо самото семантично търсене. От отговорите 76\% не съдържат грешка по същество, а 42\% са верни и изчерпателни (съдържат всички свързани с въпроса факти). Прагът за отказ сам не разпознава медицинските въпроси извън обхвата, но с проверката на ключовите понятия системата отказва на седем от десетте без нито един погрешен отказ, а нито един от десетте не получава измислено съдържание.

Качеството на отговора зависи преди всичко от извличането. Основните ограничения са малкият генеративен модел, който смесва подробности от съседни пасажи, и проблемите в данните: изкривен при сканирането текст, загубени имена на лекарства и ръководства, които дават различни дози за едно лекарство. При желание за подобряване на резултатите е добре да се ползва по-голям генеративен модел, да се фрагментира по разделите на ръководствата и да бъдат отчетени версиите им, чрез което информацията от по-новите ръководства да получи приоритет пред тази от по-старите.

\end{document}
